% Standalone document. Overleaf: choose this file as the Main document.
% Compatible with pdfLaTeX; no shell escape, custom fonts, or BibTeX required.
\documentclass[11pt,letterpaper]{article}
\usepackage[T1]{fontenc}
\usepackage[utf8]{inputenc}
\usepackage{lmodern}
\usepackage[margin=0.85in]{geometry}
\usepackage{amsmath,amssymb,booktabs,tabularx,array,microtype}
\usepackage[hidelinks]{hyperref}
\usepackage{enumitem}
\setlist{nosep,leftmargin=*}
\setlength{\parindent}{0pt}
\setlength{\parskip}{6pt}
\setlength{\emergencystretch}{2em}
\newcommand{\dd}{\mathrm d}
\title{Magnetic Levitation at a Nominal 10 mm Distance\\
\large Abridged Calculation: Position PI Driving PWM Directly}
\author{ENGR 155 calculation notes}
\date{}
\begin{document}
\maketitle
\vspace{-1.5em}

\textbf{Result.} Reducing the distance improves the available lifting force.
For the proposal's illustrative coil, a 10 mm \emph{center distance} requires
about 1.178 A and 78.5\% ideal duty at 12 V. However, assuming an inductance
of 10 mH and negligible passive mechanical damping, the position PI loop
remains unstable for every choice of PI gains in the continuous-time model
derived below. The inductance and damping are assumptions, not measurements.

\section{Geometry, architecture, and assumptions}
The suspended magnet, carrier, and payload have total mass
\(m=0.010\,\mathrm{kg}\). The desired \emph{surface gap} is
\(d_0=0.010\,\mathrm m\), measured from the lower coil face to the upper
magnet face. The thin-coil/dipole model instead uses the distance \(z\)
from the coil plane to the magnet center:
\begin{equation}
 z=d+\Delta_g,\qquad z_0=d_0+\Delta_g.
\end{equation}
For a coil represented by its axial midplane and a magnet represented by its
geometric center, \(\Delta_g\) is approximately half the coil thickness plus
half the magnet height. Its value must be determined from the actual geometry.

\textbf{Distance convention for the numbers below.} To reproduce the current
10 mm calculation, the numerical illustration sets \(z_0=10\,\mathrm{mm}\).
This equals a 10 mm surface gap only in the zero-offset, thin-coil/point-magnet
idealization. For a real 10 mm surface gap, substitute the actual
\(z_0=10\,\mathrm{mm}+\Delta_g\) into the same equations.

The controller is \emph{position PI only}: measured gap error produces PWM
duty directly. There is no current reference or current-feedback controller.
Current can still be measured for field compensation, telemetry, and protection.
The FPGA executes the position PI and PWM; the MCU handles ADC acquisition,
DMA, the gap estimate, and the web interface through the Wi-Fi module.

The numerical sizing assumptions from the proposal are
\begin{center}
\begin{tabular}{lll}
\toprule
Symbol & Meaning & Value used \\
\midrule
\(m\) & Total suspended mass & \(0.010\,\mathrm{kg}\) \\
\(g\) & Gravitational acceleration & \(9.81\,\mathrm{m/s^2}\) \\
\(N\) & Coil turns & 500 \\
\(a\) & Mean coil radius & \(0.010\,\mathrm m\) \\
\(\mu\) & Magnet dipole moment & \(0.050\,\mathrm{A\,m^2}\) \\
\(R\) & Coil resistance & \(8.0\,\Omega\) \\
\(V_s\) & Supply voltage & \(12.0\,\mathrm V\) \\
\(L\) & Differential coil inductance & \(0.010\,\mathrm H\), assumed \\
\(c\) & Passive mechanical damping & \(0\,\mathrm{N\,s/m}\), approximation \\
\bottomrule
\end{tabular}
\end{center}
The model assumes an aligned permanent dipole, a thin air-core coil, small
motions near equilibrium, and an averaged PWM voltage \(V_su\) in continuous
conduction. It neglects driver and diode drops, thermal changes, magnetic
saturation, hysteresis, and unmodeled structural motion. A final core-equipped
coil must be characterized locally rather than assigned these air-core slopes.

\section{Holding force, current, and power}
The standard axial coil field and aligned-dipole approximation give
\begin{align}
 B(z)&=\frac{\mu_0NIa^2}{2(a^2+z^2)^{3/2}},\\
 F(z,I)&\simeq \mu\left|\frac{\dd B}{\dd z}\right|=K(z)I,\\
 K(z)&=\frac{3\mu_0N\mu a^2z}{2(a^2+z^2)^{5/2}},
 \qquad \mu_0=4\pi\times10^{-7}\,\mathrm{H/m}.
\end{align}
The field expression follows the current-loop model in
OpenStax~\cite{openstax}; the force relation is the dipole approximation
used in the proposal. At \(z_0=0.010\,\mathrm m\),
\begin{align}
 K(z_0)&=0.0833041\,\mathrm{N/A},\\
 F_0=mg&=0.0981\,\mathrm N,\\
 i_0=\frac{mg}{K(z_0)}&=1.17761\,\mathrm A,\\
 u_0\simeq\frac{Ri_0}{V_s}&=0.78508,\\
 P_{\mathrm{coil}}=i_0^2R&\simeq11.09\,\mathrm W.
\end{align}
The ideal full-on current is \(V_s/R=1.50\,\mathrm A\), so this nominal
point is statically feasible in the ideal model. It leaves about 21.5
percentage points of upward duty headroom; it does not establish thermal
acceptability, dynamic force headroom, or closed-loop stability.

\section{Small-signal force gains and mechanics}
Define \(x=d-d_0\) and \(j=i-i_0\). Positive \(x\) is downward motion,
away from the coil. A first-order Taylor expansion of upward magnetic force is
\begin{equation}
 F\simeq mg-k_xx+k_fj,\qquad
 k_f=\left.\frac{\partial F}{\partial i}\right|_0,\qquad
 k_x=-\left.\frac{\partial F}{\partial d}\right|_0.
\end{equation}
Because \(z=d+\Delta_g\) with fixed geometry, \(\dd z/\dd d=1\).
For the illustrative force law,
\begin{align}
 k_f&=K(z_0),\\
 \frac{\dd K}{\dd z}
 &=K(z)\left(\frac1z-\frac{5z}{a^2+z^2}\right),\\
 k_x&=mg\left(\frac{5z_0}{a^2+z_0^2}-\frac1{z_0}\right).
\end{align}
Substitution gives
\begin{equation}
 \boxed{k_f=0.0833041\,\mathrm{N/A},\qquad k_x=14.715\,\mathrm{N/m}.}
\end{equation}
Newton's law, with downward displacement positive, is
\begin{equation}
 \boxed{m\ddot x+c\dot x-k_xx+k_fj=0.}
\end{equation}
If current were held fixed and \(c=0\), the mechanical poles would be
\(s=\pm\sqrt{k_x/m}=\pm38.36\,\mathrm{s^{-1}}\). The positive pole
represents a growing displacement. This is a fixed-current mechanical
comparison, not the complete voltage-driven plant.

\section{Position PI and the uncontrolled coil current}
Let \(u\in[0,1]\) be duty fraction. For a constant target define
\begin{equation}
 \boxed{u=u_0+K_px+K_i\eta,\qquad \dot\eta=x.}
\end{equation}
The positive sign increases attraction when the object is too far from the
coil. \(K_p\) has units \(\mathrm{m^{-1}}\), \(K_i\) has units
\(\mathrm{m^{-1}s^{-1}}\), and \(\eta\) has units \(\mathrm{m\,s}\).
For a changing requested gap \(r(t)\), the controller error is \(e=d-r\).

Coil current is a plant state, even though it has no controller. Including
motion-induced voltage, the local electrical equation is
\begin{equation}
 \boxed{L\dot j+Rj-k_f\dot x=V_s\delta u,\qquad\delta u=u-u_0.}
\end{equation}
The coupling term follows from reciprocal magnetic coenergy: the downward
magnetic force is \(-F\), so
\(\partial\lambda/\partial d=\partial(-F)/\partial i=-k_f\), where
\(\lambda\) is flux linkage. Consequently
\(\delta v=Rj+L\dot j-k_f\dot x\). This is the energy-based
electromechanical modeling method in Woodson and Melcher~\cite{woodson}.
The mechanical damping \(c\) excludes this explicitly modeled coil coupling.
If motion-induced voltage is negligible, omit the \(-k_f\dot x\) term;
then the \(+k_f^2\) term below also disappears.

\section{Laplace transform and stability calculation}
With zero initial conditions, the transformed equations are
\begin{align}
 (ms^2+cs-k_x)X+k_fJ&=0,\\
 (Ls+R)J-k_fsX&=V_s\Delta U,\\
 \Delta U&=\left(K_p+\frac{K_i}{s}\right)X.
\end{align}
Eliminating \(J\) and \(\Delta U\) gives
\begin{equation}
 s\big[(ms^2+cs-k_x)(Ls+R)+k_f^2s\big]
 +k_fV_s(K_ps+K_i)=0.
\end{equation}
Thus the characteristic polynomial is
\begin{equation}
\boxed{\begin{aligned}
 p(s)={}&mLs^4+(mR+cL)s^3+(cR-k_xL+k_f^2)s^2\\
       &+(k_fV_sK_p-k_xR)s+k_fV_sK_i.
\end{aligned}}
\end{equation}
The four states are position, velocity, current, and accumulated position
error. Taking a transform changes the representation, not the physical
stability. Laplace and feedback foundations are in Oppenheim and
Willsky~\cite{oppenheim}.

For \(p(s)=A_4s^4+A_3s^3+A_2s^2+A_1s+A_0\), strict stability requires
\begin{align}
 A_4,A_3,A_2,A_1,A_0&>0,\\
 A_3A_2-A_4A_1&>0,\\
 A_3A_2A_1-A_4A_1^2-A_3^2A_0&>0.
\end{align}
These are the quartic Routh--Hurwitz conditions~\cite{routh}. Two necessary
conditions are
\begin{equation}
 K_p>\frac{k_xR}{k_fV_s},\qquad cR+k_f^2>k_xL.
\end{equation}
The second condition is independent of the PI gains.

For \(L=10\,\mathrm{mH}\) and \(c=0\), the numerical polynomial is
\begin{equation}
\boxed{\begin{aligned}
 p(s)={}&0.0001s^4+0.08s^3-0.140210s^2\\
       &+(0.999649K_p-117.72)s+0.999649K_i.
\end{aligned}}
\end{equation}
Its \(s^2\) coefficient is negative for every \(K_p,K_i\). Therefore no
position PI tuning can make this particular continuous-time model strictly
stable. The necessary proportional-gain threshold,
\(K_p>117.76\,\mathrm{m^{-1}}\), cannot repair the negative coefficient.

For illustration only, choose \(K_p=150\,\mathrm{m^{-1}}\) and
\(K_i=100\,\mathrm{m^{-1}s^{-1}}\). Then
\begin{equation}
 p(s)=0.0001s^4+0.08s^3-0.140210s^2+32.2273s+99.9649,
\end{equation}
whose roots are
\begin{equation}
 \boxed{s=-802.247,\quad -2.9963,\quad 2.6214\pm j\,20.2237\ \mathrm{s^{-1}}.}
\end{equation}
The oscillatory mode grows as \(e^{2.6214t}\), with an amplitude
\(e\)-folding time of about 0.381 s. These are \emph{unstable example gains},
not recommended firmware settings.

The necessary electrical/damping condition can also be written
\begin{align}
 c=0:\quad L&<\frac{k_f^2}{k_x}=0.4716\,\mathrm{mH},\\
 L=10\,\mathrm{mH}:\quad c&>\frac{k_xL-k_f^2}{R}
 =0.01753\,\mathrm{N\,s/m}.
\end{align}
Passing either applicable threshold is necessary but not sufficient: all
Routh inequalities must still pass. This does not prove that every physical
PI-controlled levitator is unstable; it identifies what this assumed plant
cannot achieve.

\section{Digital implementation and the \texorpdfstring{\(z\)}{z}-transform}
For an illustrative 1 kHz position update, \(T=0.001\,\mathrm s\), choose
the explicit timing convention
\begin{align}
 \delta u[k]&=K_pe[k]+K_i\eta[k],\\
 \eta[k+1]&=\eta[k]+Te[k].
\end{align}
This gives \(C(z)=K_p+K_iT/(z-1)\). Using the updated accumulator in the
same sample instead would give a different discrete controller.

To check the held actuator correctly, define \(q=[x,\dot x,j]^\mathsf T\):
\begin{equation}
 \dot q=Aq+B\delta u,\quad
 A=\begin{bmatrix}
 0&1&0\\ k_x/m&-c/m&-k_f/m\\0&k_f/L&-R/L
 \end{bmatrix},\quad
 B=\begin{bmatrix}0\\0\\V_s/L\end{bmatrix},\quad
 C=\begin{bmatrix}1&0&0\end{bmatrix}.
\end{equation}
The exact zero-order-held plant has
\begin{equation}
 A_d=e^{AT},\qquad B_d=\int_0^T e^{A\tau}B\,\dd\tau,
\end{equation}
and the sampled closed-loop state matrix, with zero additional computation
delay, is
\begin{equation}
 M=\begin{bmatrix}A_d+B_dK_pC&B_dK_i\\TC&1\end{bmatrix}.
\end{equation}
The implemented loop is stable only if every eigenvalue of \(M\) is inside
the unit circle. Merely substituting \(z=e^{sT}\) for the continuous
closed-loop poles does not model a separately sampled controller and held
actuator. Additional sensor filtering, communication, and computation delay
must be included before selecting firmware gains.

For the same illustrative plant and the unstable example gains above,
the eigenvalues of this sampled matrix are
\begin{equation}
 z=0.447767,\quad 0.997019,\quad 1.002880\pm j\,0.020241.
\end{equation}
The complex pair has magnitude approximately 1.00308, exceeding one.
Thus the explicitly defined 1 kHz sampled implementation is also unstable.

\section{Measurements and the position-tolerance requirement}
For the final coil at the actual 10 mm surface gap, measure the holding
current and hot resistance, differential inductance, and local force slopes:
\begin{align}
 k_f&\simeq\frac{F(d_0,i_0+\Delta i)-F(d_0,i_0-\Delta i)}{2\Delta i},\\
 k_x&\simeq-\frac{F(d_0+\Delta d,i_0)-F(d_0-\Delta d,i_0)}{2\Delta d}.
\end{align}
Estimate passive damping without counting the modeled coil-induced damping
twice. Recalculate the continuous and sampled poles with the measured values.
Duty limits and integrator anti-windup are necessary practical details, but
they do not turn an unstable linear PI design into a stable small-signal loop.

Specify a true-position tolerance and a recovery time separately; for example,
\(\lvert d-r\rvert\le0.5\,\mathrm{mm}\) after a defined disturbance and
settling interval. That example is not a demonstrated performance claim.
A stable, unsaturated PI loop can reject constant disturbances in its measured
position, while sensor bias, noise, quantization, and saturation still limit
the true position. The proposal's current demonstration band is
\(\pm5\,\mathrm{mm}\). No settling time or position band can be certified
from the unstable example above.

\clearpage
\begin{thebibliography}{9}
\raggedright
\bibitem{openstax} OpenStax, \emph{University Physics, Volume 2}, Sec.~12.4,
``Magnetic Field of a Current Loop.''
\url{https://openstax.org/books/university-physics-volume-2/pages/12-4-magnetic-field-of-a-current-loop}
\bibitem{woodson} H.~H. Woodson and J.~R. Melcher,
\emph{Electromechanical Dynamics}, Part I, Chs.~3 and 5. MIT OpenCourseWare.
\url{https://ocw.mit.edu/courses/res-6-003-electromechanical-dynamics-spring-2009/pages/part-i/}
\bibitem{routh} University of Illinois, ECE 486 Handbook, Lecture 7,
``Routh--Hurwitz Stability Criterion.''
\url{https://courses.physics.illinois.edu/ece486/sp2026/documentation/handbook/lec07.html}
\bibitem{oppenheim} A.~V. Oppenheim, A.~S. Willsky, with S.~H. Nawab,
\emph{Signals and Systems}, 2nd ed., Prentice Hall, 1997,
Chs.~9 (Laplace), 10 (\(z\)-transform), and 11 (feedback).
MIT reading guide:
\url{https://ocw.mit.edu/courses/6-003-signals-and-systems-fall-2011/pages/readings/}
\end{thebibliography}
\end{document}
